% Avsnitt 15.7, ur lectures/L15/README.md.

\section{Sammanfattning}
\label{c15:sec:review}

\needspace{6\baselineskip}
\subsection*{Det här ska du kunna nu}
\begin{itemize}
  \item Implementera ett skiftregister som deserialiserar från bussen och kastar stoppbitar.
  \item Upptäcka ett stoppningsfel och säga exakt vilket villkor som utgör det.
  \item Förklara varför mottagaren samplar vid \code{sample} och sändaren skiftar vid
    \code{bit_done}.
  \item Mata CRC-motorn från mottagarsidan utan att räkna stoppbitar.
  \item Läsa någon annans VHDL och skriva en granskning som är användbar snarare än artig.
\end{itemize}

\needspace{6\baselineskip}
\subsection*{Frågor att testa dig själv med}
\begin{enumerate}
  \item Varför samplar mottagaren vid \code{sample} snarare än vid \code{bit_done}?
  \item Exakt vilket villkor utgör ett stoppningsfel, och varför är det sex lika bitar som är
    gränsen på mottagarsidan när sändaren räknar fem?
  \item Vad händer med stoppbiten när den är korrekt?
  \item Varför får \code{crc15} bara se \code{real_bit}, och vad skulle en felaktigt matad stoppbit
    ställa till?
  \item Varför räcker det inte att räkna mottagna bitar för att veta när en byte är klar?
\end{enumerate}

\needspace{6\baselineskip}
\subsection*{Kodgranskningsseminariet}
Ungefär en timme, i helgrupp. Varje grupp får tio minuter och visar en mergad Pull Request (vad den
ändrade, hur den var uppdelad, och vad commit-meddelandena säger), en granskning gruppen skrivit
(vad granskaren tittade efter, och vad som faktiskt ändrades av kommentarerna), och en sak som inte
fungerat i arbetssättet, och vad ni gjort åt det.

Välj ut i gruppen en mergad Pull Request och en granskning någon i gruppen skrivit att visa. Det
behöver inte vara den finaste; en PR som fick befogad kritik är mer värd att visa än en som ingen
läste.

Seminariet är obligatoriskt och är den ena av projektets två fasta hållpunkter. Det betygsätts inte
för sig, men det ingår i den helhetsbedömning av arbetssättet som beskrivs i
bilaga~\ref{app:project}.

\needspace{6\baselineskip}
\subsection*{Riktvärde}
Med \code{rx_shift_reg} klar finns alla fyra lövmoduler, och \code{can_controller} är den enda som
återstår på CAN-sidan. Grupper som ligger efter bör lägga tiden här snarare än att börja på
tillståndsmaskinen: den behöver alla fyra.

\needspace{6\baselineskip}
\subsection*{Blicken framåt}
\Chapref{c16:ch} tar ramen som tillståndsmaskin, först förklarad och sedan byggd på sändarsidan.
